% Generated draft: requires mathematical and source review.
\documentclass[11pt]{article}
\usepackage[margin=25mm]{geometry}
\usepackage{fontspec}
\setmainfont{Latin Modern Roman}
\usepackage{amsmath,amssymb,mathtools,unicode-math}
\setmathfont{Latin Modern Math}
\usepackage{enumitem,needspace,xurl,hyperref}
\setlength{\parindent}{0pt}
\setlength{\parskip}{5pt}
\setlength{\emergencystretch}{3em}
\newcommand{\N}{\mathbb{N}}
\newcommand{\Z}{\mathbb{Z}}
\newcommand{\Q}{\mathbb{Q}}
\newcommand{\R}{\mathbb{R}}
\newcommand{\C}{\mathbb{C}}
\newcommand{\F}{\mathbb{F}}
\newcommand{\A}{\mathbb{A}}
\newcommand{\PP}{\mathbb{P}}
\newcommand{\E}{\mathbb{E}}
\newcommand{\eps}{\varepsilon}
\newcommand{\dd}{\,\mathrm{d}}
\begin{document}
\section*{UnsolvedMath ID: 20001652}
\textbf{Source problem code:} catalogue entry
\subsection{Definitions and mathematical statement}
Problem 4.3 asks whether every homomorphism \(\varphi:\operatorname{Out}(F_n)\to\operatorname{MCG}(S)\) has virtually abelian image. A group is virtually abelian if it has an abelian finite-index subgroup; every finite group qualifies, without needing to be abelian. The source does not specify rank, orientability, punctures, boundary or mapping-class conventions. The aid's quantitative theorem fixes \(n\ge3\) and \(\operatorname{Mod}(S_g)\), the orientation-preserving mapping class group of a closed orientable surface with \(g\ge2\).

Let
\[
W_n=(\Z/2\Z)^n\rtimes S_n,\qquad |W_n|=2^n n!,
\]
the signed-permutation subgroup of \(\operatorname{Out}(F_n)\), obtained by inverting and permuting a free basis. The aid reports that any infinite-image homomorphism to \(\operatorname{Mod}(S_g)\) forces the target to contain an actual isomorphic copy of \(W_n\). It attributes this implication to the Bridson--Vogtmann alternative combined with Farb--Masur superrigidity. Since finite subgroups of the closed orientation-preserving genus-\(g\) mapping class group have order at most \(84(g-1)\), an infinite image requires
\[
2^n n!\le84(g-1).
\]
Consequently every image is finite whenever
\[
84(g-1)<2^n n!.
\]
Under the stronger inequality \(84(g-1)<(n+1)!\), the aid reports the sharper conclusion \(|\varphi(\operatorname{Out}(F_n))|\le2\). Equality at either threshold is not included in these strict finite-image ranges.

For \(n\ge4\), the aid invokes Property (T) for \(\operatorname{Out}(F_n)\) to identify virtually abelian image with finite image: quotients inherit Property (T), and a virtually abelian group with that property is finite. This reduction is not asserted for \(n=3\) from the same argument, nor do the genus bounds prove the general all-genus question. For an orientable surface with nonempty boundary fixed pointwise, the aid's same package gives triviality of homomorphisms from the full \(\operatorname{Out}(F_n)\) in its stated \(n\ge3\) setting. Pointwise-fixed boundary conventions are not interchangeable with punctures or boundary permutations.

The surrounding source discussion is materially different from an exact numerical theorem. It tentatively attributes to ``Martin Bridson, I think'' the statement that the genus should be at least \(n!\), based on rose symmetries. The supplied verification says that \(n!\) genus assertion really is printed, rather than an OCR error. It also clarifies that the damaged expression ``Z2 o S n'' is the signed-permutation wreath product and the numerical expression is \(2^n n!\), not a group \(\Z^n\) or an equality between a group and a number. The draft retains the printed genus claim as contextual source wording and separately states the aid's precise \(84(g-1)\) inequality; they are not silently identified as the same bound.

The source further says that ``Poison groups'' in \(\operatorname{Out}(F_n)\) have virtually abelian image by results of ``Martin and Juan,'' but does not define those groups or provide a theorem reference. It explicitly warns that finite-index subgroups change the story. The signed finite-subgroup argument above is a result for the full source group; it is not automatically a theorem for an arbitrary torsion-free finite-index subgroup, which may not contain the required finite symmetry group.

The registry labels the entry partially solved; attempt 1 is a partial result, problem open at the run, internal rigor proved and novelty confidence low. The aid is machine generated, not peer reviewed, and requires independent verification before citation.
\subsection{Short English statement}
The question asks whether maps from \(\operatorname{Out}(F_n)\) to mapping class groups have virtually abelian image. In the aid's closed orientable-surface setting, an infinite image requires a signed-permutation subgroup large enough to force \(2^n n!\le84(g-1)\); smaller genera give finite images. The original printed genus remark and finite-index warning are retained separately, without claiming a general all-genus solution.
\subsection{Sources}
\begin{itemize}
\item Local UnsolvedMath registry, permanent ID \(20001652\), source code \texttt{catalogue entry}; \texttt{aim-geometric-group-theory-notes.json}, \texttt{notes[160]}.
\item AIM, \emph{The geometry of the outer automorphism group of a free group}, Problem 4.3 and discussion with tentative Martin Bridson attribution and ``Martin and Juan'' remark: \url{https://aimath.org/WWN/outerauto/outerauto.pdf}.
\item Bridson--Vogtmann alternative, Farb--Masur superrigidity, finite-subgroup surface bound and modern Property (T) results, as named/described by the supplied research summary. No full titles, years or theorem numbers are provided. The original literature field is empty.
\item Supplied research attempt 1 and registry update August 14, 2026. The verification says the \(n!\) genus remark is printed in the PDF; no independent PDF inspection is claimed.
\end{itemize}
\end{document}
